\subsection{Translation fields exclude general-type components}

\begin{statement}{C.2.UN.15}{Proved}{No component of a torus degeneration has general-type normalization}
Let $f:U\to\Delta$ be a proper holomorphic map from
a smooth complex threefold to a disc, with connected
fibres and smooth fibres complex two-tori over
$\Delta^*$. For every irreducible component $S_i$
of the reduced central support, its smooth
normalization resolution $Z_i$ satisfies
\[
                            \kappa(Z_i)\le1.
 \tag{C.SC.1}
\]
No K\"ahler, projectivity, Betti-number or Euler-number
assumption is imposed on $U$ or its central fibre.
\end{statement}
\begin{proof}
We first make the chosen component reduced by an
actual finite base change. If its full-fibre
multiplicity is $m_i$, put $t=s^{m_i}$ and normalize
the fibre product. At a generic smooth point of
$S_i$, away from the other components, the original
equation is $t=u z^{m_i}$ with $u$ a unit. The
normalized base change has the explicit local maps
\[
                  w^{m_i}=u,\qquad z=s/w.
 \tag{C.SC.2}
\]
It is therefore smooth at those points, and every
central component above $S_i$ has multiplicity one
and maps generically finitely onto $S_i$. Choose one
such component and resolve the normalized total
space and central support, isomorphically over
$\Delta^*$ and generically along that component.
Write $f':U'\to\Delta_s$ for the resulting proper
map and $E$ for the smooth strict component. The
space $U'$ is smooth, $E$ still has multiplicity
one, and the smooth nearby fibres are the original
connected tori. Connectedness of all full fibres
follows also from Stein factorization over the
normal disc, since its generic degree is one.

If $Z_i$ were of general type, so would be $E$.
Here is the precise comparison. The map $E\to S_i$
factors through normalization and gives a
generically finite meromorphic map $E\dashrightarrow
Z_i$. Resolve its graph to obtain a birational
map $a:\widehat E\to E$ and a generically finite
holomorphic map $b:\widehat E\to Z_i$ between
smooth compact surfaces. The actual Jacobian map
gives
\[
                 K_{\widehat E}=b^*K_{Z_i}+R,
                 \qquad R\ge0.                 \tag{C.SC.3}
\]
Pullback of pluricanonical sections is injective.
The generically finite pullback of a big canonical
line is big; equivalently a generically finite
pluricanonical map of $Z_i$ remains generically
finite after composition with $b$. Equation
(C.SC.3) therefore makes $K_{\widehat E}$ big.
Birational invariance gives $\kappa(E)=2$.
In particular $E$ is projective, as is every
smooth compact complex surface of general type.

We construct a nonzero meromorphic vector field
on $E$ with single-valued solutions. Define the
coherent relative tangent sheaf by the actual
kernel
\[
 \mathcal T_{U'/\Delta_s}
 =\ker\bigl[df':T_{U'}\longrightarrow f'^*T_{\Delta_s}\bigr].
 \tag{C.SC.4}
\]
Its proper direct image is torsion-free on the
disc: a section annihilated by a power of $s$
vanishes on the dense open punctured tube and
hence vanishes identically. Its rank is two,
because the global holomorphic vector fields
of a two-torus are its two translation fields.
It is thus a locally free rank-two sheaf after
shrinking the disc. Choose a nonzero section
$v$, viewed by adjunction as an actual
holomorphic vertical vector field on $U'$.
At the generic point of $E$ its order of
vanishing is a finite integer $k\ge0$.
Set
\[
                              w=s^{-k}v.
 \tag{C.SC.5}
\]
This is a global meromorphic vector field on
$U'$, with poles confined to the central
divisor. At a generic point of $E$, the
multiplicity-one assertion makes $s$ a local
transverse coordinate. Since $v(s)=0$, write
$v=a(x,y,s)\partial_x+b(x,y,s)\partial_y$.
The integer $k$ is the minimum of the two
coefficient orders along $s=0$. Thus $w$
is holomorphic and nonzero generically on
$E$ and has a well-defined nonzero tangent
meromorphic restriction $w_E$ to $E$.

We use the precise closed-univalence theorem
\cite[Theorem~3.3, with its proof]{Guillot2024}:
for a holomorphic vector field on a complex
manifold of any dimension, the locus where
its maximal solutions are single-valued is
closed. Such a field is called semicomplete
when that locus is the whole manifold.
The closure assertion is a path-continuation
statement. A failure at a point is witnessed
by a nonclosed path in a solution leaf whose
time-form integral is zero. Lifting the
time parametrization of this compact path
through finitely many flow charts produces
the same nonclosed zero-time path through
all sufficiently nearby initial points.
This proves openness of the failure locus
and records exactly the comparison used here.

Let $P$ be the polar divisor of $w$ and put
$W=U'\setminus P$. On the dense open subset
$U'^*\subset W$, $w$ restricts to a
holomorphic vector field on each compact
torus, multiplied by the constant $s^{-k}$
on that fibre. Each such field is complete:
its complex-time solution stays in the
compact fibre, so local existence extends
it for all times. Consequently $w$ is
semicomplete on all of $W$ by the stated
closedness theorem.

The hypersurface $E\cap W$ is invariant.
Indeed $w(s)=0$, and at any point where
$w$ is holomorphic, factoring the local
equation of $s$ into its irreducible
components shows that $w$ preserves the
ideal of each component. A solution
starting in $E\cap W$ therefore stays
there, and its single-valuedness is the
one inherited from $w$. The curve removed
from $E$ by the ambient polar divisor may
contain removable poles of $w_E$. Apply
the same closedness theorem once more on
$E$ minus the actual polar divisor of
$w_E$. Its dense open subset $E\cap W$
is already semicomplete. Hence $w_E$ is
a nonzero semicomplete meromorphic vector
field on the entire compact surface $E$,
in the precise sense of single-valuedness
on its holomorphic locus.

We finish by proving that a projective
surface of general type admits no such
field. Suppose first that one maximal
solution of $w_E$ has Zariski-dense image.
The Main Theorem of \cite{Guillot2019},
including the clarification immediately
following it, gives a birational model
which is an abelian surface or admits a
rational or elliptic fibration. The
holomorphic alternative is covered there
by Proposition~22; the dense leaf excludes
a first integral in that alternative.
All these models have Kodaira dimension
at most one, contradicting $\kappa(E)=2$.
Only single-valued maximality, already
proved, is needed in this application;
the solution need not be entire.

Here is the other case, with its actual
fibration map. If every solution leaf is
contained in a proper analytic curve,
the closures of the regular leaves are
irreducible algebraic invariant curves
on the projective surface $E$. They
cover its open regular foliation locus.
Use the countably many Hilbert polynomials
and flattening strata of their universal
families to parametrize reduced irreducible
curves. The locus of such integral curves
is constructible. After clearing the
polar divisor, invariance is the vanishing
of the induced normal-component map of
the vector field on the universal curve;
on these strata it is a constructible
condition. Thus we have countably many
irreducible locally closed parameter
strata of integral invariant curves.
Their incidence images in $E$ are
constructible. If none were dense, their
Zariski closures would be countably many
proper closed analytic subsets covering
the regular foliation locus, contrary
to the Baire theorem. Choose a stratum
whose incidence dominates $E$. Its
generic curves are reduced and
irreducible. At a regular point of the
foliation exactly one such irreducible
curve can pass: two would contain the
same local solution germ and would
coincide. Thus the parameter has dimension
one and the incidence map to $E$ is
generically one-to-one. Compactify and
normalize this parameter curve and
close the incidence graph. Inverting
its birational projection to $E$ and
projecting to the parameter gives an
actual meromorphic first integral.
Resolve its indeterminacies and take its
Stein factor to obtain a morphism
\[
                         q:E'\longrightarrow C
 \tag{C.SC.6}
\]
with connected generic fibre a smooth
compact curve. The lifted meromorphic
field remains semicomplete: it agrees
with the old field on the birational
open set, and the same closedness theorem
extends single-valuedness over its full
holomorphic locus. It is a nonzero
tangent meromorphic field on a generic
fibre $G$ of (C.SC.6).

Such a field on a compact curve has no
pole. A pole of order $r\ge1$ would
have the local form $z^{-r}a(z)\partial_z$
with $a(0)\ne0$. Its time form is
$z^r a(z)^{-1}dz$, whose primitive has
a zero of order $r+1$. The inverse
solution therefore has $r+1$ branches
around that time value, all lying in
the punctured coordinate disc. This
contradicts semicompleteness. The field
on $G$ is consequently holomorphic
and nonzero, so
$\deg T_G=2-2g(G)\ge0$ and $g(G)\le1$.

This also contradicts general type
directly through pluricanonical maps.
Adjunction on the smooth generic fibre
gives $K_{E'}|_G=K_G$, of degree at
most zero. Every linear system
$|nK_{E'}|$ therefore restricts to a
constant map on $G$: a negative-degree
line has no sections, and a degree-zero
line on a connected compact curve has
at most one independent section. A
generically finite pluricanonical map
of the general-type surface $E'$ cannot
contract its generic fibres in this
way. This excludes the all-algebraic-
leaves case as well.

Both cases contradict the general-type
consequence of (C.SC.3). The chosen
$Z_i$ cannot have Kodaira dimension
two. Since $S_i$ was arbitrary,
(C.SC.1) follows for every component.
\end{proof}
